\newpage
\lesson{17/11}
Consider now
\[
    A=\bra{0}T(\psi(x_1)\psi(x_2)\bar\psi(x_3)\bar\psi(x_4))\ket{0}
\]
one can use Wick's theorem, with the \textit{caveat} of putting the contracted
fields next to one another (and gaining - signs):
\[
    A=-\overline{\psi_1\bar\psi_3}\,\overline{\psi_2\bar\psi_4} +
    \overline{\psi_1\bar\psi_4}\,\overline{\psi_2\bar\psi_3}.
\]
With the Path Integral, this is
\[
    A = \eval{\frac{\delta^4}{\delta\bar\eta_1\delta\bar\eta_2\delta\eta_3
    \delta\eta_4}Z_0[\eta,\bar\eta]}_{\eta_i,\bar\eta_j = 0}
\]
after verifying the sign is correct, one can use
\[
    Z_0[\eta,\bar\eta] = \exp\left[ -\left\langle \bar\eta(x)S(x-y)\eta(y)
    \right\rangle_{xy}\right]
\]
to find the vacuum expectation value:
    \begin{align*}
        &A=\eval{\frac{\delta^4}{\delta\bar\eta_1\delta\bar\eta_2\delta\eta_3
        \delta\eta_4}\left\{\exp\left[ -\left\langle \bar\eta(x)S(x-y)\eta(y)
        \right\rangle_{xy}\right]\right\}}_{\eta,\bar\eta=0}=
        \eval{\frac{\delta^3}{\delta\bar\eta_1\delta\bar\eta_2\delta\eta_3}
        \left\langle \bar\eta(x)S_{xx_4}\right\rangle_x \exp\left[ ... \right]
        }_{\eta,\bar\eta=0}=\\
        &=\eval{\frac{\delta^2}{\delta\bar\eta_1\delta\bar\eta_2}
        \left(-\left\langle \bar\eta(x)S_{xx_4}\right\rangle_x\left\langle
        \bar\eta(y)S_{yx_3}\right\rangle_y\exp\left[ ... \right]\right)}_{
        \eta,\bar\eta=0}=
        \eval{\fdv[]{}{\bar\eta_1}\left(-S_{24}\left\langle \bar\eta(y)S_{yx_3}
        \right\rangle_y+\left\langle \bar\eta(x)S_{xx_4}\right\rangle_xS_{23}
        \right)}_{\eta,\bar\eta=0}=\\
        & =-S_{24}S_{13}+S_{14}S_{23}
    \end{align*}
    which is the same as Wick's theorem. Once again, path integrals automatically
    apply Wick's theorem to amplitudes' computations.
    
\section{Interacting fermionic theory: Yukawa}
An interaction term like
\[
    \int_{}^{} \dd^4x a_\mu\pd^\mu\psi
\]
would break Lorentz invariance (why? maybe spinor transformations).\\
Odd number of fields would go to zero and the simplest interaction term
(Fermi theory) has problems with renormalization (that will be clear in the
Advanced QFT course).

Another field needs to be introduced, a scalar field $\phi$ of mass $M$:
\[
    S = \int_{}^{} \dd^4x \bar\psi(i\slashed{\pd}-m)\psi+\underbrace{
    \frac{1}{2}\pd_\mu\phi\pd^\mu\phi - \frac{M^2}{2}\phi^2}_\text{pure scalar}-
    \underbrace{g\phi\bar\psi\psi - \alpha(\bar\psi\psi)}_\text{interactions}
\]
where $g,\alpha$ are coupling constants and $m,M$ are \textit{bare} masses of the
fields.

Since $\left[\phi\right]=L^{-1}=M$ and $\left[\psi\right] = M^{\flatfrac{3}{2}}$,
$\left[g\right]=1$ (which is good) and $\left[\alpha\right]=M^{-2}$: since
$\alpha$ is dimensional, it must be multiplied by $E^2$ to become dimensionless,
but this is counterproductive for perturbative approaches, since at higher
energies $\alpha E^2$ grows; the term has to be discarded.
\[
    S = \int_{}^{} \dd^4x \bar\psi(i\slashed{\pd}-m)\psi+\underbrace{
    \frac{1}{2}\pd_\mu\phi\pd^\mu\phi - \frac{M^2}{2}\phi^2}_\text{pure scalar}-
    \underbrace{g\phi\bar\psi\psi}_\text{interaction}
\]
This can be written as
\[
    S = S_0 - g\int_{}^{} \dd^4x \phi\bar\psi\psi
\]
For the corresponding free theory, the propagator is
\[
    Z_0[\eta,\bar\eta,J] = \frac{1}{\mathcal{N}} \int \mathcal{D}\phi\mathcal{D}
    \psi\mathcal{D}\bar\psi \exp\left[ iS_0 + i\left\langle J\phi\right\rangle
    +i\left\langle\bar\eta\psi\right\rangle+i\left\langle\bar\psi\eta
    \right\rangle \right]=\frac{1}{\mathcal{N}}
    \exp\left[ \frac{-i}{2}\left\langle J(x)\Delta_{xy}J(y) \right\rangle_{xy}-
    \left\langle \bar\eta(x)S_{xy}\eta(y)\right\rangle_{xy}\right]
\]
The interacting propagator uses $S$ instead of $S_0$ and for small $g$ it can be
formally Taylor-expanded
\[
    Z[J,\eta,\bar\eta] = \sum_{n=0}^{\infty} \frac{(-ig)^n}{n!}\int \mathcal{D}\phi
    \mathcal{D}\psi\mathcal{D}\bar\psi
    (\left\langle\phi\bar\psi\psi\right\rangle_x)^n e^{iS_0+...}
\]
and the fields can be brought down by deriving it
\[
    \frac{\delta^3}{\delta J\delta\eta\delta\bar\eta}\int \mathcal{D}\phi
    \mathcal{D}\psi\mathcal{D}\bar\psi\exp\left[ iS_0 + i\left\langle J\phi
    \right\rangle+i\left\langle \bar\eta\psi\right\rangle +i\left\langle \bar\psi
    \eta\right\rangle\right] 
\]  
this brings down $3$ $i$ factors and $2$ - signs and another $-$ comes from the
anticommutation of $\psi\bar\psi = - \bar\psi\psi$, so to get the
vacuum expectation value a $-i$ factor has to be multiplied to the derivative.

So the full expression of the propagator becomes
\[
    Z[J,\eta,\bar\eta] = \exp\left[ -ig\left\langle \frac{-i\delta^3}{\delta J
    \delta\eta\delta\bar\eta}\right\rangle_x\right]Z_0[J,\eta,\bar\eta]=
    \left(1-g \left\langle \frac{\delta^3}{\delta J\delta\eta\delta\bar\eta}
    \right\rangle_x+... \right)Z_0[J,\eta,\bar\eta].
\]
The free energy $W[J,\eta,\bar\eta]$ is $-\ln Z$: this can be treated
exactly like the $\phi^4$ interaction and at first order it is
\begin{gather*}
    W[J,\eta,\bar\eta] = W_0[J,\eta,\bar\eta] - g e^{W_0}\left\langle 
    \frac{\delta^3}{\delta J\delta\eta\delta\bar\eta}\right\rangle_xe^{-W_0}
    +O(g^2)
\end{gather*}
the $O(g)$ term is
\begin{gather*}
    e^{W_0}\left\langle \frac{\delta^3}{\delta J\delta\eta\delta\bar\eta}
    \right\rangle_xe^{-W_0}=e^{W_0}\left\langle \frac{\delta^2}{\delta J\delta\eta
    }(-1)\left\langle S_{xy}\eta_y\right\rangle_y\right\rangle_xe^{-W_0}=\\=
    e^{W_0}\left\langle \fdv[]{}{J}\left(S_{xx}-\left\langle S_{xy}\eta_y
    \right\rangle_y \left\langle \bar\eta_yS_{yx}\right\rangle_y\right)
    e^{-W_0}\right\rangle_x=
    \left\langle \left(S_{xx}-\left\langle S_{xy}\eta_y\bar\eta_uS_{ux}
    \right\rangle_{yu} \right) \Delta_{xw}J_w \right\rangle_{xw}
\end{gather*}
The full free energy is then
\[
    W[J,\eta,\bar\eta] = \underbrace{-\frac{1}{2}\left\langle J_x\Delta_{xy}J_y
    \right\rangle_{xy} - \left\langle \bar\eta_xS_{xy}\eta_y\right\rangle_{xy}}_
    {=W_0}-g
    \left(
    \left\langle S_{xx}\Delta_{xy}J_y\right\rangle_{xy} +
    \left\langle\bar\eta_uS_{ux}S_{xy}\eta_y\Delta_{xw}J_w\right\rangle_{xywu}
    \right) 
\]
Since $\eta$ are spinors, the indexes must be considered, since
\[
    \fdv[]{\eta_\alpha(x)}{\eta_\beta(y)}=\delta_{\alpha\beta}\delta^4(x-y)
\]
this is implemented in the last expression, since $\bar\eta$ are column vectors,
$\eta$ are row vectors, $S$ are matrices and $\Delta,J$ are scalars.
Note that the first $S_{xx}$ is a trace, i.e. has spinor indices $\alpha\alpha$,
so it is a scalar.
